\documentclass[11pt,a4paper]{article}
\usepackage[utf8]{inputenc}
\usepackage[T1]{fontenc}
\usepackage[margin=2.5cm]{geometry}
\usepackage{amsmath,amssymb}
\usepackage{booktabs}
\usepackage{hyperref}

\usepackage{array}

\newcommand{\tx}[1]{\texttt{#1}}

\DeclareMathOperator{\lcm}{lcm}

\title{NSUCRYPTO 2026, Round 2\\Problem 5: ``Grandchildren''}
\date{}
\author{Jacob Lu, Raymond Ouyang, Ruiyang Ye}
\begin{document}
\maketitle

\section*{Answer}

\[
\text{Asya } 4,\qquad \text{Vanya } 6,\qquad \text{Sonya } 12 \qquad(\text{key } 4\text{-}6\text{-}12),
\]
with house number $H=192$ and Grandma at $G=46$, on the same (even) side of a street with odd/even numbering.
The only ingredients beyond the text are a model of Sonya's walk and of the street (Section~\ref{sec:cond}); no
assumption about the children's ages is added. The answer is the only candidate that survives every one of the
24 model settings we tested.

\section{Formalisation}

Let $a$, $v$, $s$ be the ages of Asya, Vanya and Sonya (positive integers), $H$ their house number, and
$G=\overline{a\,v}$ Grandma's house number. Sonya passes $N=6s$ houses on the way. The value of $N$ depends on
how the street is numbered:
\begin{itemize}
  \item \textbf{Consecutive numbering} ($1,2,3,\dots$ along the walk): $N=|G-H|-1$.
  \item \textbf{Odd/even numbering} (both houses on the same side, as is usual on Russian streets):
  $G\equiv H\pmod 2$ and $N=\tfrac12|G-H|-1$.
\end{itemize}

\section{All conditions}\label{sec:cond}

\begin{center}
\begin{tabular}{@{}cll@{}}
\toprule
& Condition & Source \\
\midrule
A & $a\mid H$, $v\mid H$, $s\mid H$, $G\neq H$; Sonya passes $6s$ houses (either numbering) & text \\
B & Vanya has math homework, so he is of school age: $v\ge 6$ & text \\
C & Sonya's 25-minute walk gives a realistic plot width (see below) & text + model \\
D & The street has a realistic length (see below) & model \\
E & The three are children: $1\le a,v,s\le 17$ & text \\
F & The key (ages of Asya, Vanya, Sonya written in order) has at least 4 digits & lock safety \\
\bottomrule
\end{tabular}
\end{center}

\paragraph{Conditions C and D.} Sonya walks for $1500$\,s. We assume typical walking speeds by age: $0.90$\,m/s
up to age~4, $1.05$ at 5--6, $1.15$ at 7--9, $1.25$ at 10--12, and $1.35$\,m/s from~13. She passes $6s$ houses,
so one plot is
\[
w(s)=\frac{\text{speed}(s)\cdot 1500\ \text{s}}{6s+1}
\]
wide. C asks that $w(s)$ lie in a realistic range. The plot width also turns house numbers into distances: the
farthest house lies about $\max(G,H)$ plots from the start of the street under consecutive numbering, or
$\max(G,H)/2$ plots under odd/even numbering. D asks that this length stay below a cap.

\paragraph{Condition F.} The ages form the digital key to Grandma's house. A 3-digit code (only 1000 combinations)
is not safe for a house lock, so we exclude every triple in which all three ages are single digits.

\paragraph{Model settings.} C and D depend on modelling parameters, so we evaluate them in 24 settings: speeds
scaled by $\times0.8$, $\times1.0$ or $\times1.2$; plot widths of 8--40\,m or 6--50\,m; and street caps of 3, 5,
8 or 12\,km. A candidate's score is the number of settings it survives, and the answer is the candidate that
survives the most.

\paragraph{Method.} We enumerated every $(a,v,s)$ with ages $1$--$17$ and every $H\le10^5$; since
$|G-H|\le2(6\cdot17+1)$ and $G\le1717$, this covers all solutions. Condition~A alone has 127 solutions with
consecutive numbering and 170 with odd/even numbering. We then applied B--F to all of them in one pass (script
\tx{one\_shot.py}).

\section{Result}

Forty candidates survive at least one setting. Those surviving at least 18 of the 24:

\begin{center}
\begin{tabular}{@{}clrrrrr@{}}
\toprule
$(a,v,s)$ & numbering & $H$ & $G$ & $w(s)$ & street & settings \\
\midrule
$(4,6,12)$   & odd/even    & 192 & 46  & 25.7\,m & 2.5\,km & \textbf{24/24} \\
$(2,7,14)$   & consecutive & 112 & 27  & 23.8\,m & 2.7\,km & 22/24 \\
$(2,9,10)$   & consecutive & 90  & 29  & 30.7\,m & 2.8\,km & 22/24 \\
$(4,6,16)$   & odd/even    & 240 & 46  & 20.9\,m & 2.5\,km & 22/24 \\
$(5,8,15)$   & odd/even    & 240 & 58  & 22.3\,m & 2.7\,km & 22/24 \\
$(1,11,11)$  & consecutive & 44  & 111 & 28.0\,m & 3.1\,km & 20/24 \\
$(1,10,14)$  & odd/even    & 280 & 110 & 23.8\,m & 3.3\,km & 20/24 \\
$(1,15,9)$   & odd/even    & 225 & 115 & 31.4\,m & 3.5\,km & 20/24 \\
$(1,17,17)$  & odd/even    & 323 & 117 & 19.7\,m & 3.2\,km & 20/24 \\
$(11,8,12)$  & odd/even    & 264 & 118 & 25.7\,m & 3.4\,km & 20/24 \\
$(3,14,12)$  & odd/even    & 168 & 314 & 25.7\,m & 4.0\,km & 18/24 \\
$(16,6,14)$  & odd/even    & 336 & 166 & 23.8\,m & 4.0\,km & 18/24 \\
\bottomrule
\end{tabular}
\end{center}

\noindent ($w(s)$ and the street length are given at speed $\times1.0$.) In the central setting (speed
$\times1.0$, plots 8--40\,m, cap 5\,km) the twelve candidates above all survive; the model settings are what
separate them. Only $(4,6,12)$ survives all 24.

\paragraph{The answer in detail.} For $(a,v,s)=(4,6,12)$ there is exactly one solution: $H=192=16\cdot\lcm(4,6,12)$,
divisible by 4, 6 and 12; $G=\overline{4\,6}=46$; both are even, so the houses are on the same side; and Sonya
passes $\tfrac12(192-46)-1=72=6\cdot12$ houses. Vanya is of school age (B), and the key 4-6-12 has four digits (F).
At age~12 Sonya walks at 1.25\,m/s, so one plot is $1.25\cdot1500/73$\,m wide, and the farthest house is
$192/2=96$ plots away:

\begin{center}
\begin{tabular}{@{}crrcc@{}}
\toprule
speed & $w(12)$ & street & plots 8--40\,m, caps 3/5/8/12\,km & plots 6--50\,m, caps 3/5/8/12\,km \\
\midrule
$\times0.8$ & 20.5\,m & 1.97\,km & \checkmark\ \checkmark\ \checkmark\ \checkmark & \checkmark\ \checkmark\ \checkmark\ \checkmark \\
$\times1.0$ & 25.7\,m & 2.47\,km & \checkmark\ \checkmark\ \checkmark\ \checkmark & \checkmark\ \checkmark\ \checkmark\ \checkmark \\
$\times1.2$ & 30.8\,m & 2.96\,km & \checkmark\ \checkmark\ \checkmark\ \checkmark & \checkmark\ \checkmark\ \checkmark\ \checkmark \\
\bottomrule
\end{tabular}
\end{center}

\noindent Its plot widths lie in the middle of the allowed range, and even its longest street is below the
strictest cap of 3\,km. 

\section{What each condition contributes}

For each specific condition dropped:

\begin{center}
\begin{tabular}{@{}lrl@{}}
\toprule
Dropped & survivors & top score \\
\midrule
none            & 40 & $(4,6,12)$ alone at 24/24 \\
F (key length)  & 53 & $(4,6,12)$ ties with $(1,6,9)$ (key 1-6-9) at 24/24 \\
B (homework)    & 96 & 14 candidates tie at 24/24 \\
C (plot width)  & 47 & $(4,6,12)$ ties with $(1,14,7)$ at 24/24 \\
D (street)      & 58 & 27 candidates tie at 24/24 \\
C and D         & 85 & all 85 tie: no model, no answer \\
\bottomrule
\end{tabular}
\end{center}
This shows that every condition is neccessary here to pinpoint a good candidate in our model.
\section*{Conclusion}

The digital key is \textbf{4-6-12}: Asya~4, Vanya~6, Sonya~12, living at number 192, with Grandma's house number 46.
The only modelling ingredients we assumed are the speed
table, the plot widths and the street caps, giving the unique candidate that passes our model.
\end{document}
